\begin{table}[!htbp]
    \centering
    \caption{Attrition balance: extensive specification}
    \label{tab:attrition_extensive}
    \scriptsize
    \begin{adjustbox}{width=\textwidth}
    \begin{tabular}{lccccc}
    \toprule
     & \multicolumn{2}{c}{Batch 1} & \multicolumn{3}{c}{Batch 2} \\
    \cmidrule(lr){2-3} \cmidrule(lr){4-6}
     & Blocked/Unavailable & Any Attrition & Not Yet Scraped & Blocked/Unavailable & Any Attrition \\
    \midrule
    Treated SMI & 0.004* & 0.008* & -0.037* & -0.008 & -0.044 \\
     & (0.002) & (0.005) & (0.020) & (0.009) & (0.027) \\
    Pure control mean & 0.005 & 0.013 & 0.149 & 0.095 & 0.244 \\
    Observations & 51542 & 51542 & 31393 & 31393 & 31393 \\
    Current extensive sample restrictions & Yes & Yes & Yes & Yes & Yes \\
    Country and batch FE & Yes & Yes & Yes & Yes & Yes \\
    Permutation-based SE & Yes & Yes & Yes & Yes & Yes \\
    \bottomrule
    \end{tabular}
    \end{adjustbox}
    \floatfoot{\textit{Notes:} Columns report the effect of being assigned to one treated initially-followed SMI on endline attrition outcomes for the current extensive-analysis sample. The sample applies the same one-influencer, baseline-posting, and percentile restrictions used in the main extensive regressions. Batch 1 and Batch 2 are shown separately; the not-yet-scraped outcome is only relevant for Batch 2 because Batch 1 scraping is complete. Standard errors are permutation-based standard deviations computed from 1,000 reassigned treatment draws.}
\end{table}
